\documentclass[11pt]{article}
\usepackage[utf8]{inputenc}
\usepackage{graphicx}
\usepackage{amsmath, amssymb}
\usepackage{geometry}
\geometry{a4paper, margin=1in}
\usepackage{hyperref}
\usepackage{caption}

\title{Parameterizing Pathological Gait with Limit Cycle Oscillators: \\ From Pure Phase to Full Amplitude Resilience}
\author{}
\date{\today}

\begin{document}
\maketitle

\begin{abstract}
Continuous biological motions like human walking represent stable limit cycles, yet distinguishing the subtle structural instabilities of pathological gait (e.g., following a stroke) from natural variations due to walking speed remains a challenge. In this note, we compile two specific phenomenological modeling approaches applied to empirical high-dimensional marker trajectories: an initial phase-only oscillator model (based on the Kuramoto/Shaevitz framework) and a full amplitude-phase Stuart-Landau model. By characterizing both the global cycle stability and the localized step-to-step volatility of spatial parameters, we uncover differences between Able-Bodied (AB) individuals and Hemiparetic subjects.
\end{abstract}

\section{Phase-Only Oscillators: The Kuramoto Framework}
Initially, we approached the cyclic motion via a purely phase-based framework. Human marker coordinates were embedded into a delay-coordinate Hankel matrix and projected via Singular Value Decomposition (SVD) to identify the two primary dominant harmonic cycles. The phase angles of these modes were fit to a coupled oscillator framework to evaluate the temporal coupling between joints.

Figure~\ref{fig:shaevitz_fit} demonstrates the model tracking phase progression. By exploring the parameters driving this timing constraint, we constructed a UMAP projection of the individual human trials (Figure~\ref{fig:shaevitz_umap}). While the model accurately parameterizes the timing characteristics of the limit cycle, subsequent analysis of the ``phase resilience"---a global timing-recovery metric---revealed strong confounding correlations with the subject's self-selected walking speed (Figures~\ref{fig:shaevitz_resilience} and~\ref{fig:shaevitz_residuals}). This suggested that while stroke pathways disrupt gait, modeling pure timing excludes critical information about spatial instability and Floquet multiplier structure.

\begin{figure}[h!]
    \centering
    \begin{minipage}{0.48\textwidth}
        \centering
        \includegraphics[width=\linewidth]{../Shaevitz_model/shaevitz_fit.png}
        \caption{Coupled phase oscillator fit.}
        \label{fig:shaevitz_fit}
    \end{minipage}\hfill
    \begin{minipage}{0.48\textwidth}
        \centering
        \includegraphics[width=\linewidth]{../Shaevitz_model/umap_shaevitz_parameters.png}
        \caption{UMAP clustering of the phase-only parameter regime across clinical populations.}
        \label{fig:shaevitz_umap}
    \end{minipage}
    
    \vspace{0.5cm}
    \begin{minipage}{0.48\textwidth}
        \centering
        \includegraphics[width=\linewidth]{../Shaevitz_model/phase_resilience_out.png}
        \caption{Phase resilience tightly tracks simple walking speed rather than pathology alone.}
        \label{fig:shaevitz_resilience}
    \end{minipage}\hfill
    \begin{minipage}{0.48\textwidth}
        \centering
        \includegraphics[width=\linewidth]{../Shaevitz_model/resilience_residuals.png}
        \caption{Residuals of phase resilience upon factoring out speed.}
        \label{fig:shaevitz_residuals}
    \end{minipage}
\end{figure}

\section{Encoding Amplitude: The Stuart-Landau Model}
To capture structural spatial limit-cycle phenomena, we upgraded to a Stuart-Landau (Hopf) network, tracking the complex amplitudes $z(t)$:
\[ \frac{dz_i}{dt} = \mu_i z_i - z_i|z_i|^2 + i \omega_i z_i + \sum_j K_{ij} z_j \]
A central benefit of this model is the direct accessibility of the structural dampening parameter $\mu$, reflecting Floquet-like decay toward the attractor. Fitting dynamic experimental gradients without unbounded derivative blowups required enforcing strong biologically constrained empirical frequencies ($\omega \approx \omega_{actual}$) and strict Ridge regularization on coupling magnitudes $K_{ij}$ evaluated via central differences. The globally fitted parameters correctly integrate autonomously to replicate the fundamental harmonics of observed human gait modes (Figure~\ref{fig:sl_val}).

\begin{figure}[h!]
    \centering
    \includegraphics[width=0.75\textwidth]{../Stuart_Landau_model/stuart_landau_validation_fixed_freq.png}
    \caption{Autonomous forward simulation of the regularized Stuart-Landau fit successfully tracking biological multi-channel amplitude variations over time.}
    \label{fig:sl_val}
\end{figure}

\subsection{Global Resilience Correlates}
Computing the global structural resilience metric ($\mu$) confirmed strong dependence on velocity. Utilizing a non-parametric localized LOESS regression exclusively on Able-Bodied users, we mapped the baseline ``healthy" limit-cycle rigidity across kinematic walking speeds. 

\begin{figure}[h!]
    \centering
    \begin{minipage}{0.48\textwidth}
        \centering
        \includegraphics[width=\linewidth]{../Stuart_Landau_model/sl_resilience_with_loess.png}
        \caption{Spatial structural rigidity $\mu$ vs. Velocity. Note the sliding-window baseline characterizing healthy subjects (AB).}
        \label{fig:sl_res}
    \end{minipage}\hfill
    \begin{minipage}{0.48\textwidth}
        \centering
        \includegraphics[width=\linewidth]{../Stuart_Landau_model/sl_mu_residuals_loess.png}
        \caption{Residual structural resilience of Stroke groups relative to the healthy non-parametric baseline.}
        \label{fig:sl_res_residuals}
    \end{minipage}
\end{figure}


\section{Intra-Trial Volatility: The Step-by-Step Variance}
While the global parameters highlighted a systematic deficit in stroke patients (reflecting a clinical resilience delay of roughly $\sim$20ms), such a minor average margin might not capture the true mechanism of pathological locomotor instability. We hypothesized the deficit operates dynamically: stroke patients frequently shift attractor stabilization strategies step-by-step to compensate for motor noise, contrasting with the uniform, fixed-rigidity cycles of healthy pedestrians.

By executing the Stuart-Landau inversion across binned temporal windows (e.g., estimating local $\mu$ parameters using 300 temporally consecutive tracking frames per step), we quantified intra-trial deviations. Rather than evaluating standard residuals, we quantified the inherent variation of $\mu$ over time. We plotted the intra-trial statistical volatility (Standard Deviation) of the stability corrections over time against local velocities (Figure \ref{fig:binned_vel}) and distributed by condition (Figure \ref{fig:binned_box}). 

\begin{figure}[h!]
    \centering
    \begin{minipage}{0.48\textwidth}
        \centering
        \includegraphics[width=\linewidth]{../Stuart_Landau_model/binned_mu_vs_speed_loess.png}
        \caption{Binned spatial rigidity evaluation over single-step local temporal sliding windows.}
        \label{fig:binned_vel}
    \end{minipage}\hfill
    \begin{minipage}{0.48\textwidth}
        \centering
        \includegraphics[width=\linewidth]{../Stuart_Landau_model/binned_delta_mu_std_boxplot.png}
        \caption{Intra-trial Volatility of Spatial Resilience across populations. Pathological patients display significantly rougher dynamic recalibration from step to step.}
        \label{fig:binned_box}
    \end{minipage}
\end{figure}

\end{document}